\chapter{The Formalism of Quantum Mechanics}
\label{ch:formalism}

The wave mechanics of the previous chapters can be recast in the language of vector spaces: a state is a vector, an observable a Hermitian operator, and a wavefunction just the components of the state in one particular basis. This chapter motivates the formalism with spin, develops the linear algebra, and restates the postulates. (The basic definitions of vector spaces, bases and linear maps are in the mathematical-methods notes.)

\section{The Stern--Gerlach Experiment}
% ============================================================

\begin{figure}[ht]
\centering
\begin{tikzpicture}[>=Stealth]
  \draw (0,0) ellipse (0.6 and 0.4); \node at (0,0) {\small Ag oven};
  \draw[thick] (1.6,0.15) -- (1.6,0.8); \draw[thick] (1.6,-0.15) -- (1.6,-0.8);
  \node[below] at (1.6,-0.85) {\small collimator};
  \draw[thick] (3.2,0.5) -- (3.5,0.9) -- (4.3,0.9) -- (4.6,0.5) -- (4.2,0.5) -- (3.9,0.2) -- (3.6,0.5) -- cycle; \node at (3.9,0.7) {S};
  \draw[thick] (3.2,-0.5) rectangle (4.6,-0.9); \node at (3.9,-0.7) {N};
  \node[below] at (3.9,-0.95) {\small magnet, $\partial B_z/\partial z\neq0$};
  \draw[thick] (6.5,-1.2) -- (6.5,1.2); \node[right] at (6.5,0) {\small screen};
  \draw[green!50!black] (0.6,0) -- (3.2,0) -- (6.5,0.8);
  \draw[green!50!black] (3.2,0) -- (6.5,-0.8);
\end{tikzpicture}
\caption{The Stern--Gerlach apparatus. A beam of silver atoms passes through an inhomogeneous field and splits into two.}
\end{figure}

Because $\partial B_z/\partial z\neq0$, the field exerts a force on objects with a magnetic moment $\bm{\mu}$:
\[
    \Delta H = -\bm{\mu}\cdot\vect{B}
    \quad\Longrightarrow\quad
    \vect{F} = -\grad\Delta H = \mu_z\pdv{B_z}{z}\,\khat .
\]
For a classical current loop of an electron orbiting a nucleus (charge $Q$, radius $R$, mass $M$), $I = \lambda v$ with $\lambda = Q/2\pi R$ and $|\vect{L}| = MRv$, so $|\bm{\mu}| = I\pi R^2 = \frac{Q}{2}Rv = \frac{Q}{2M}|\vect{L}|$. With no preferred orientation, one expects a continuous spread $-\mu_{\max}\le\mu_z\le\mu_{\max}$.

\begin{fact}[What is observed]
Instead, $\mu_z = \pm\mu$, with equal frequency. Fundamental particles such as the electron have an intrinsic magnetic moment
\[
    \bm{\mu} = g\mu_e\vect{J}/\hbar,
    \qquad
    \mu_e = \frac{e\hbar}{2m_e},
    \qquad
    g = 2.00231930436092(36),
\]
with $g$ the gyromagnetic ratio ($g = 2$ from relativity, with quantum corrections). So the experiment measures $J_z = S_z = \pm\hbar/2$. Orbital angular momentum must be an integer multiple of $\hbar$, so this is a new intrinsic property: \emph{spin}. (Silver has 47 electrons: the filled shells have zero net spin, and the last electron is in a $5s$ orbital with no orbital angular momentum. The nucleus has a very small moment, so $\mu_{\text{Ag}}\approx\mu_e$.)
\end{fact}

The apparatus ``measures $S_z$'', and we can chain such devices:
\begin{enumerate}
    \item $\to$ SG($\hat{z}$) $\to$ $S_z = \pm\hbar/2$, 50\% each. Blocking $S_z = -\hbar/2$ and passing the rest through a second SG($\hat{z}$) gives $S_z = +\hbar/2$ 100\% of the time.
    \item $\to$ SG($\hat{x}$) gives $S_x = \pm\hbar/2$, 50\% each.
    \item SG($\hat{z}$), keep $+$, then SG($\hat{x}$): $S_x = \pm\hbar/2$, 25\% each of the original beam.
    \item SG($\hat{z}$) keep $+$, SG($\hat{x}$) keep $+$, then SG($\hat{z}$): $S_z = \pm\hbar/2$, 12.5\% each. A state with both $S_z = +\hbar/2$ and $S_x = +\hbar/2$ does not survive: the $S_x$ measurement erased the $S_z$ information.
\end{enumerate}

% ============================================================
\section{Spin-\texorpdfstring{$\tfrac12$}{1/2} States}
\label{sec:spin-half-intro}
% ============================================================

These experiments need a new formalism. After SG($\hat{z}$) the electron is in one of two states,
\[
    |\hat{z},+\rangle\leftrightarrow\begin{pmatrix}1\\0\end{pmatrix},
    \qquad
    |\hat{z},-\rangle\leftrightarrow\begin{pmatrix}0\\1\end{pmatrix},
\]
and by experiment 1 the initial state must be $|\psi\rangle = c_1|\hat{z},+\rangle + c_2|\hat{z},-\rangle = (c_1,c_2)^T$ for some $c_i\in\mathbb{C}$. Measuring $S_z$ is implemented by the matrix
\[
    \hat{S}_z\leftrightarrow\frac\hbar2\begin{pmatrix}1&0\\0&-1\end{pmatrix},
    \qquad
    \hat{S}_z|\hat{z},\pm\rangle = \pm\frac\hbar2|\hat{z},\pm\rangle .
\]
Similarly, after SG($\hat{x}$) the state is $|\hat{x},+\rangle$ or $|\hat{x},-\rangle$, and the same initial state can be written $|\psi\rangle = \tilde{c}_1|\hat{x},+\rangle + \tilde{c}_2|\hat{x},-\rangle$ (experiment 2). Experiment 4 shows that measuring $\hat{S}_x$ changes the $\hat{S}_z$ state: the operations do not commute. Rotations, generated by orbital angular momentum, also fail to commute, $[\hat{L}_i,\hat{L}_j] = i\hbar\epsilon_{ijk}\hat{L}_k$, so we try the same algebra for spin.

\begin{formula}[Spin-$\tfrac12$ operators]
\[
    \hat{S}_x = \frac\hbar2\begin{pmatrix}0&1\\1&0\end{pmatrix},
    \qquad
    \hat{S}_y = \frac\hbar2\begin{pmatrix}0&-i\\i&0\end{pmatrix},
    \qquad
    \hat{S}_z = \frac\hbar2\begin{pmatrix}1&0\\0&-1\end{pmatrix},
\]
i.e.\ $\hat{S}_i = \frac\hbar2\sigma_i$, satisfying $[\hat{S}_i,\hat{S}_j] = i\hbar\epsilon_{ijk}\hat{S}_k$.
\end{formula}

\begin{fact}[Pauli matrices]
$\sigma_i^\dagger = \sigma_i$, $\operatorname{tr}\sigma_i = 0$, $\sigma_i^2 = I$, and $[\sigma_i,\sigma_j] = 2i\epsilon_{ijk}\sigma_k$.
\end{fact}

The eigenstates of $\hat{S}_x$ follow at once:
\[
    |\hat{x},\pm\rangle = \frac{1}{\sqrt2}|\hat{z},+\rangle \pm \frac{1}{\sqrt2}|\hat{z},-\rangle
    \leftrightarrow\frac{1}{\sqrt2}\begin{pmatrix}1\\\pm1\end{pmatrix}.
\]
Each basis can be expressed in terms of the other, consistent with the intensities seen in experiments 3 and 4.

% ============================================================
\section{Inner Products and Hilbert Space}
% ============================================================

\begin{definition}[Inner product]
An inner product on a vector space $V$ over $\mathbb{F} = \mathbb{R}$ or $\mathbb{C}$ assigns to each ordered pair $u,v\in V$ a number $\langle u,v\rangle\in\mathbb{F}$ such that
\begin{enumerate}
    \item $\langle u,v\rangle = \langle v,u\rangle^*$ (conjugate symmetry);
    \item $\langle u,v_1+v_2\rangle = \langle u,v_1\rangle + \langle u,v_2\rangle$ and $\langle u,av\rangle = a\langle u,v\rangle$ (linearity in the second argument);
    \item $\langle v,v\rangle\ge0$, with equality if and only if $v = 0$.
\end{enumerate}
Properties 1 and 2 imply conjugate linearity in the first argument, $\langle au,v\rangle = a^*\langle u,v\rangle$. The norm is $|v| = \sqrt{\langle v,v\rangle}$.
\end{definition}

Vectors are orthogonal if $\langle u,v\rangle = 0$, in which case $|u+v|^2 = |u|^2+|v|^2$ (Pythagoras).

\begin{theorem}[Schwarz and triangle inequalities]
\[
    |\langle u,v\rangle|\le|u|\,|v|,
    \qquad
    |u+v|\le|u|+|v|,
\]
the first saturated if and only if $v = cu$ for some $c\in\mathbb{F}$.
\end{theorem}
\begin{proof}
Let $w = u - \frac{\langle v,u\rangle}{|v|^2}v$. Then $0\le|w|^2 = |u|^2 - |\langle u,v\rangle|^2/|v|^2$.
\end{proof}

\begin{example}[Inner products]
On $\mathbb{C}^2$, $\langle a,b\rangle = a_1^*b_1 + a_2^*b_2$. On functions $[0,L]\to\mathbb{C}$, $\langle f,g\rangle = \int_0^Lf^*g\,dx$; for instance $\sin(\pi x/L)$ and $\sin(2\pi x/L)$ are orthogonal.
\end{example}

An \emph{orthonormal basis} has $\langle e_i,e_j\rangle = \delta_{ij}$, and then $v = \sum_ia_ie_i$ with $a_i = \langle e_i,v\rangle$. Any basis $\{v_i\}$ can be made orthonormal by the \emph{Gram--Schmidt procedure}: $e_1 = v_1/|v_1|$, then $\tilde{v}_j = v_j - \sum_{i<j}\langle e_i,v_j\rangle e_i$ and $e_j = \tilde{v}_j/|\tilde{v}_j|$. For a subspace $U$, the orthogonal complement $U^\perp = \{v : \langle v,u\rangle = 0\ \forall u\in U\}$ gives $V = U\oplus U^\perp$, and the orthogonal projection onto $U$ is $P_Uv = \sum_i\langle e_i,v\rangle e_i$ over an orthonormal basis of $U$.

\begin{definition}[Hilbert space]
A complex vector space with an inner product is a \emph{Hilbert space} $\mathcal{H}$. If it is infinite-dimensional we also require completeness: every Cauchy sequence converges to a vector in $\mathcal{H}$.
\end{definition}

% ============================================================
\section{Dirac Notation}
% ============================================================

Following Dirac, vectors are written as \emph{kets} $|v\rangle$, with $|u+v\rangle = |u\rangle+|v\rangle$ and $|av\rangle = a|v\rangle$, and the inner product as a \emph{bra-ket} $\langle u|v\rangle$. The \emph{bra} $\langle u|$ is not a vector in $V$: it is the linear functional $V\to\mathbb{C}$ that sends $|v\rangle\mapsto\langle u|v\rangle$, an element of the dual space $V^*$.

In an orthonormal basis, a ket is a column vector and its bra the conjugate row vector:
\[
    |v\rangle = \sum_iv_i|e_i\rangle\leftrightarrow\begin{pmatrix}v_1\\v_2\\\vdots\end{pmatrix},
    \qquad
    \langle v| = \sum_iv_i^*\langle e_i|\leftrightarrow(v_1^*,v_2^*,\ldots),
    \qquad
    \langle v|v\rangle = \sum_i|v_i|^2 .
\]
For functions, $\langle f|$ is the instruction ``$\int_0^Lf^*(x)\,\underline{\quad}\,dx$''.

% ============================================================
\section{Linear Operators}
% ============================================================

A linear operator $\hat{T}:V\to V$ satisfies $\hat{T}(|u\rangle+|v\rangle) = \hat{T}|u\rangle + \hat{T}|v\rangle$ and $\hat{T}(a|v\rangle) = a\hat{T}|v\rangle$. On polynomials, multiplication by $3x^2-5x+2$, differentiation, and $\int_{-1}^1dx\,(2x+3)p(x)$ are linear; cubing is not, and dividing by $x$ leaves the space. The set $\mathcal{L}(V)$ of linear operators is itself a vector space, and with the product $(ST)|v\rangle = S(T|v\rangle)$ an algebra. Multiplication is associative but not commutative, and the commutator $[S,T] = ST-TS$ is again an operator.

The \emph{null space} $\{|v\rangle : T|v\rangle = 0\}$ and \emph{range} $\{T|v\rangle\}$ are subspaces with $\dim V = \dim\operatorname{null}T + \dim\operatorname{range}T$. $T$ is injective iff $\operatorname{null}T = \{0\}$, and invertible if some $S$ has $ST = TS = I$.

\begin{formula}[Matrix elements]
In a basis $\{|v_i\rangle\}$, $T$ is fixed by its action on the basis vectors,
\[
    T|v_j\rangle = \sum_iT_{ij}|v_i\rangle,
    \qquad
    (TS)_{kj} = \sum_iT_{ki}S_{ij},
\]
so the matrix of a product is the product of the matrices. If $A|v_j\rangle = |u_j\rangle$ relates two bases, then $T(\{u\}) = A^{-1}T(\{v\})A$.
\end{formula}

\begin{definition}[Adjoint]
The adjoint $T^\dagger$ is defined by $\langle u|Tv\rangle = \langle T^\dagger u|v\rangle$ for all $u,v$. Then $(ST)^\dagger = T^\dagger S^\dagger$, $(S^\dagger)^\dagger = S$, and in an orthonormal basis $(T^\dagger)_{ij} = T_{ji}^*$: the conjugate transpose. Acting on bras, $\langle u|T = \langle T^\dagger u|$.
\end{definition}

\begin{definition}[Hermitian, unitary and normal operators]
\leavevmode
\begin{itemize}
    \item $T$ is \emph{Hermitian} (self-adjoint) if $T^\dagger = T$, equivalently $\langle v|T|v\rangle\in\mathbb{R}$ for all $v$.
    \item $U$ is \emph{unitary} if it is surjective and preserves norms, $|U|v\rangle| = ||v\rangle|$. Then $U^\dagger U = UU^\dagger = I$, so $U^\dagger = U^{-1}$, and $U$ preserves inner products, $\langle Uu|Uv\rangle = \langle u|v\rangle$. It maps an orthonormal basis to an orthonormal basis.
    \item $T$ is \emph{normal} if $[T,T^\dagger] = 0$. Hermitian, anti-Hermitian ($T^\dagger = -T$) and unitary operators are normal.
\end{itemize}
\end{definition}

\begin{formula}[Outer products and projectors]
Operators can be built from outer products, $T = \sum_k|a_k\rangle\langle b_k|$, acting as $T|v\rangle = \sum_k\langle b_k|v\rangle|a_k\rangle$; in a basis $|a\rangle\langle b|$ is the matrix $a_ib_j^*$. The projector onto $|m\rangle$ is $P_m = |m\rangle\langle m|$, with $P_m^2 = P_m$. For an orthonormal basis,
\[
    I = \sum_i|e_i\rangle\langle e_i|,
\]
the \emph{resolution of the identity}. A unitary relating two orthonormal bases, $|f_i\rangle = U|e_i\rangle$, is $U = \sum_i|f_i\rangle\langle e_i|$.
\end{formula}

% ============================================================
\section{Eigenvalues and the Spectral Theorem}
% ============================================================

A subspace $U$ is $T$-invariant if $T|u\rangle\in U$ for all $|u\rangle\in U$. A one-dimensional invariant subspace gives $T|u\rangle = \lambda|u\rangle$: $\lambda$ is an \emph{eigenvalue} and $|u\rangle$ an \emph{eigenvector}. The eigenvectors with eigenvalue $\lambda$ form $\operatorname{null}(T-\lambda I)$, so $\lambda$ solves the characteristic equation $\det(T-\lambda I) = 0$. The set of eigenvalues is the \emph{spectrum}; an eigenvalue with several independent eigenvectors is \emph{degenerate}, with that number its multiplicity. Eigenvectors with distinct eigenvalues are linearly independent. Eigenvalues are basis-independent, and $\det T = \prod\lambda_i$, $\operatorname{tr}T = \sum\lambda_i$.

\begin{theorem}[Hermitian operators]
The eigenvalues of a Hermitian operator are real, and eigenvectors with distinct eigenvalues are orthogonal.
\end{theorem}
\begin{proof}
$\lambda\langle v|v\rangle = \langle v|Tv\rangle = \langle Tv|v\rangle = \lambda^*\langle v|v\rangle$. For $T|v_i\rangle = \lambda_i|v_i\rangle$, $\lambda_1\langle v_2|v_1\rangle = \langle v_2|Tv_1\rangle = \langle Tv_2|v_1\rangle = \lambda_2\langle v_2|v_1\rangle$.
\end{proof}

$T$ is \emph{diagonalizable} if some basis consists of its eigenvectors, so that $T_{ij} = \lambda_i\delta_{ij}$; not every operator is (e.g.\ $\begin{pmatrix}0&1\\0&0\end{pmatrix}$). It is \emph{unitarily diagonalizable} if that basis can be orthonormal, so that $U^\dagger TU = D$ for a unitary $U$ whose columns are the eigenvectors.

\begin{theorem}[Spectral theorem]
On a complex vector space, the eigenvectors of $T$ form an orthonormal basis if and only if $T$ is normal. In particular every Hermitian operator is unitarily diagonalizable. Grouping the eigenvalues $\lambda_k$ with degeneracies $d_k$, $V = U_1\oplus U_2\oplus\cdots$ with $U_k$ the eigenspaces, and
\[
    T = \sum_k\lambda_kP_k,
    \qquad
    P_k = \sum_{i=1}^{d_k}|u_i^{(k)}\rangle\langle u_i^{(k)}|,
    \qquad
    P_kP_l = \delta_{kl}P_k,
    \qquad
    \sum_kP_k = I .
\]
\end{theorem}

\begin{theorem}[Simultaneous diagonalization]
Commuting Hermitian operators can be simultaneously diagonalized, and so can any set of mutually commuting operators.
\end{theorem}

In quantum mechanics, commuting observables are \emph{compatible}: each can be measured without disturbing a previously measured value of the other. If a Hermitian operator is non-degenerate, measuring it uniquely fixes the state. If it is degenerate, the state is only known to lie in the eigenspace.

\begin{definition}[Complete set of commuting observables]
A CSCO is a set of mutually commuting operators $\{S_1,\ldots,S_n\}$ whose eigenvalues uniquely label every joint eigenstate. Given an incomplete set, add operators that commute with all of them but split the remaining degeneracies, until it is complete. CSCOs are not unique.
\end{definition}

\begin{example}[A CSCO]
With $A = \operatorname{diag}(1,1,2)$, $B = \operatorname{diag}(1,2,2)$ and $C = \operatorname{diag}(1,2,3)$, all commuting: neither $\{A\}$ nor $\{B\}$ is a CSCO, but $\{A,B\}$ and $\{C\}$ are, labelling the basis as $|a_i,b_i\rangle$ or $|c_i\rangle$.
\end{example}

% ============================================================
\section{Continuous Bases}
% ============================================================

For a particle on a line, introduce position eigenstates $\hat{x}|x\rangle = x|x\rangle$ labelled by a continuous variable. (The label is not the vector: $|x_0+x_1\rangle\neq|x_0\rangle+|x_1\rangle$.)

\begin{formula}[Position and momentum bases]
\[
    \langle x|y\rangle = \delta(x-y),
    \qquad
    I = \int dx\,|x\rangle\langle x|,
    \qquad
    |\psi\rangle = \int dx\,\psi(x)|x\rangle,
    \qquad
    \psi(x) = \langle x|\psi\rangle,
\]
\[
    \langle p|p'\rangle = \delta(p-p'),
    \qquad
    I = \int dp\,|p\rangle\langle p|,
    \qquad
    \langle x|p\rangle = \frac{e^{ipx/\hbar}}{\sqrt{2\pi\hbar}},
    \qquad
    \tilde\psi(p) = \langle p|\psi\rangle .
\]
\end{formula}

The wavefunction is the component of $|\psi\rangle$ along $|x\rangle$. Overlaps become $\langle\varphi|\psi\rangle = \int dx\,\langle\varphi|x\rangle\langle x|\psi\rangle = \int\varphi^*\psi\,dx$, and the Fourier transform is just a change of basis, $\langle p|\psi\rangle = \int dx\,\langle p|x\rangle\langle x|\psi\rangle$. The states $|x\rangle$ and $|p\rangle$ are not normalizable, so strictly they are not in the Hilbert space; including them (a ``rigged Hilbert space'') is the pragmatic approach. Inserting complete sets:
\[
    \langle x|\hat{p}|\psi\rangle = \frac\hbar i\frac{d}{dx}\langle x|\psi\rangle,
    \qquad
    \langle p|\hat{x}|\psi\rangle = i\hbar\frac{d}{dp}\langle p|\psi\rangle,
    \qquad
    [\hat{x},\hat{p}] = i\hbar\,\hat{I}.
\]

% ============================================================
\section{The Postulates}
% ============================================================

\begin{fact}[Postulates of quantum mechanics]
\leavevmode
\begin{enumerate}
    \item The state of a physical system is a vector $|\Psi\rangle$ in a Hilbert space $\mathcal{H}$.
    \item Every measurable quantity is described by a Hermitian operator on $\mathcal{H}$.
    \item The only possible outcomes of a measurement of $\hat{A}$ are its eigenvalues. Measuring $\hat{A}$ on a normalized $|\Psi\rangle$ gives $\alpha_n$ with probability $P(\alpha_n) = |\langle\alpha_n|\Psi\rangle|^2$, after which the state is $|\alpha_n\rangle$ (not $\hat{A}|\Psi\rangle$).
    \item Time evolution is unitary: $|\Psi(t)\rangle = \hat{U}(t,t_0)|\Psi(t_0)\rangle$. This implies the Schr\"odinger equation $i\hbar\frac{d}{dt}|\Psi\rangle = \hat{H}|\Psi\rangle$ for a Hermitian $\hat{H}$ (Chapter~\ref{ch:dynamics}).
\end{enumerate}
\end{fact}

% ============================================================
\section{Uncertainty Relations}
% ============================================================

The expectation value $\langle A\rangle_\Psi = \langle\Psi|A|\Psi\rangle = \sum_\alpha\alpha P(\alpha)$ is real for Hermitian $A$.

\begin{definition}[Uncertainty]
\[
    \Delta A(\Psi) = \bigl|(\hat{A}-\langle A\rangle I)|\Psi\rangle\bigr|,
    \qquad
    \Delta A^2 = \langle A^2\rangle - \langle A\rangle^2 .
\]
$\Delta A = 0$ if and only if $|\Psi\rangle$ is an eigenstate of $\hat{A}$. Geometrically, $\Delta A$ is the length of the component of $\hat{A}|\Psi\rangle$ perpendicular to $|\Psi\rangle$, i.e.\ of $(1-|\Psi\rangle\langle\Psi|)\hat{A}|\Psi\rangle$.
\end{definition}

\begin{theorem}[Uncertainty relation]
For Hermitian $A$, $B$ and a normalized state,
\[
    \Delta A^2\,\Delta B^2\ge\Bigl(\Bigl\langle\Psi\Bigl|\frac{1}{2i}[A,B]\Bigr|\Psi\Bigr\rangle\Bigr)^2 .
\]
\end{theorem}
\begin{proof}
Let $|f\rangle = (A-\langle A\rangle)|\Psi\rangle$ and $|g\rangle = (B-\langle B\rangle)|\Psi\rangle$, so $\Delta A^2 = \langle f|f\rangle$ and $\Delta B^2 = \langle g|g\rangle$. By Schwarz, $\Delta A^2\Delta B^2\ge|\langle f|g\rangle|^2 = (\mathrm{Re}\langle f|g\rangle)^2 + (\mathrm{Im}\langle f|g\rangle)^2$, and $\mathrm{Im}\langle f|g\rangle = \frac{1}{2i}(\langle f|g\rangle - \langle g|f\rangle) = \langle\Psi|\frac{1}{2i}[A,B]|\Psi\rangle$. Drop the real part. (The operator $\frac{1}{2i}[A,B]$ is Hermitian.)
\end{proof}

The bound is saturated when $|g\rangle = i\lambda|f\rangle$ with $\lambda\in\mathbb{R}$, so that $\mathrm{Re}\langle f|g\rangle = 0$; then $|\lambda| = \Delta B/\Delta A$. For $[\hat{x},\hat{p}] = i\hbar$ this gives $\Delta x\,\Delta p\ge\hbar/2$. If $[A,B] = 0$ the relation is empty, and compatible observables can be known exactly together.

\begin{example}[Bounding a ground-state energy]
For $\hat{H} = \hat{p}^2/2m + \alpha\hat{x}^4$, the ground state can be chosen real, so $\langle\hat{p}\rangle = 0$ and $\langle\hat{p}^2\rangle = \Delta p^2$. Also $\langle x^4\rangle\ge\langle x^2\rangle^2\ge\Delta x^4$. With $\Delta p\ge\hbar/2\Delta x$,
\[
    \langle H\rangle_0\ge\frac{\hbar^2}{8m\Delta x^2} + \alpha\Delta x^4\ge\min_{\Delta x}(\ldots) = 2^{1/3}\frac38\Bigl(\frac{\hbar^2\sqrt\alpha}{m}\Bigr)^{2/3}\approx0.47\Bigl(\frac{\hbar^2\sqrt\alpha}{m}\Bigr)^{2/3},
\]
compared with the numerical value $0.67(\hbar^2\sqrt\alpha/m)^{2/3}$.
\end{example}
